\clearpage
\thispagestyle{empty}
\begin{rpgdark}[TABELA DE ROTAS · d20]
\centering
{\textgoth{\fontsize{22}{24}\selectfont TABELA DE ROTAS}}\\[0.25em]
{\scshape\bfseries\small VINTE PORTAS PARA O ARQUIVO}\\[0.3em]
{\color{rpggold}\rule{0.5\textwidth}{0.6pt}}
\end{rpgdark}

\vspace{0.4em}
\noindent{\small\itshape Role 1d20 e escolha uma entrada. O texto não muda; a rota só decide por onde você entra. As páginas são as desta edição.}

\vspace{0.25em}
\begin{rpgbox}[ENTRADAS]
\footnotesize
\renewcommand{\arraystretch}{0.98}
\begin{tabularx}{\textwidth}{@{}c l >{\raggedright\arraybackslash}X r@{}}
\toprule
\textbf{d20} & \textbf{Fragmento} & \textbf{Título} & \textbf{Página} \\
\midrule
1 & MEM-01 & O Sertão Antes do Fim & \pageref{frag:MEM-01} \\
2 & MEM-04 & A Vala & \pageref{frag:MEM-04} \\
3 & MEM-10 & O Primeiro Homicídio & \pageref{frag:MEM-10} \\
4 & MEM-14 & O Diário Emerge & \pageref{frag:MEM-14} \\
5 & MEM-20 & A Morte do Vizinho de Cela & \pageref{frag:MEM-20} \\
6 & MEM-25 & A Morte no Espelho & \pageref{frag:MEM-25} \\
7 & MEM-27 & A Página Seguinte & \pageref{frag:MEM-27} \\
8 & DOC-01 & A Exumação & \pageref{frag:DOC-01} \\
9 & DOC-03 & Laudo de Necropsia & \pageref{frag:DOC-03} \\
10 & DOC-06 & Relatório de Sessão III & \pageref{frag:DOC-06} \\
11 & DOC-13 & Diário de Oliveira III & \pageref{frag:DOC-13} \\
12 & DOC-19 & Holocausto Brasileiro & \pageref{frag:DOC-19} \\
13 & DOC-20 & Termo de Cadeia de Custódia & \pageref{frag:DOC-20} \\
14 & LUC-01 & O Convite & \pageref{frag:LUC-01} \\
15 & LUC-02 & O Arquivo & \pageref{frag:LUC-02} \\
16 & LUC-09 & A Exumação de Oliveira & \pageref{frag:LUC-09} \\
17 & LUC-17 & O Medo & \pageref{frag:LUC-17} \\
18 & LUC-20 & O Arquivo Aberto & \pageref{frag:LUC-20} \\
19 & CONT-01 & Você Abriu o Diário & \pageref{frag:CONT-01} \\
20 & CONT-13 & A Fila & \pageref{frag:CONT-13} \\
\bottomrule
\end{tabularx}
\end{rpgbox}

\vspace{0.4em}
\noindent{\footnotesize\itshape A tabela é um aparato de leitura; não altera a ordem, o texto ou o desfecho. O índice completo está nas páginas de navegação.}
\clearpage
